\documentclass[10pt]{article}
\usepackage[margin=20mm]{geometry}
\usepackage{amsmath,amssymb,mathtools,fontspec,unicode-math,xurl,hyperref}
\IfFontExistsTF{DejaVu Serif}{\setmainfont{DejaVu Serif}}{\setmainfont{Latin Modern Roman}}
\setmathfont{latinmodern-math.otf}
\setlength{\parindent}{0pt}\setlength{\parskip}{6pt}\sloppy
\title{Technical verification of abelian varieties
not isogenous to Jacobians over global fields}
\author{AI-assisted technical review}
\date{}
\begin{document}
\maketitle
Abstract: The review follows the finite-jet and monodromy construction, the Serre–Tate calculation, the moving-divisor reduction, and the intersection argument over number and finite fields. Two local defects are established: reduction modulo four alone does not force full \(2\)-adic general linear monodromy, and inverse \(p\)-multiplication in multiplicative Serre–Tate coordinates requires inverse Frobenius on coefficients. The paper’s additional finite-quotient condition repairs the first inference; a uniform calculation on reduced formal branches shows why the second does not invalidate the required jet-preservation conclusion. Three further items require clarification or additional proof: the meaning of group constituents, reduced images versus correspondence cycles, and the Galois condition in Lemma 6.2. A precise reduced-image counterexample is given without claiming to refute the valid cycle interpretation. The highest established impact is Minor; unresolved items do not raise that assessment. Scientific adjudication remains pending.\par
\section*{Identity, scope and quantifiers}
The source is A. N. Shankar and J. Tsimerman, Abelian varieties not isogenous to Jacobians over global fields, arXiv:2105.02998v3, 5 March 2025, twenty-two pages [paper]. The frozen journal identity is Duke Mathematical Journal 174 (2025), no. 18, 3535–3559, DOI https://doi.org/10.1215/00127094-2025-0016. The acquired twenty-two-page author version is labelled final by its arXiv metadata, but the twenty-five-page journal text was not acquired; their equivalence is not assumed. PDF and printed page numbers agree.\par
All six sections, the bibliography, and the actual source TeX were read. Material following `\textbackslash{}end\{document\}` and commented alternative arguments were distinguished from the article. Theorems 1.1 and 1.2 concern \(F=\overline{\mathbb F_p(t)}\), the algebraic closure of the rational function field. Theorem 1.1 permits a divisor over \(F\), whereas Theorem 1.2 first treats a divisor defined over a finite field. Both assume \(g>1\). The Jacobian application uses the properness of the Torelli locus for \(g\geq4\); it is not a conclusion about every abelian variety of every dimension. Theorem 1.3 concerns a Hodge-generic point over \(\overline{\mathbb Q}\). The finite-field assertion is \[\exists x\in\mathcal A_g(\overline{\mathbb F}_p)^{\rm ord}\quad
  \forall \tau\text{ prime to }p,\qquad \tau(x)\not\subset D .\] It does not say that the entire isogeny orbit is disjoint from \(D\).\par
The auxiliary odd prime \(r\ne p\) rigidifies the moduli problem. Maximal monodromy is understood relative to the group allowed by that fixed level. In characteristic zero its finite index in the full adelic similitude group must be absorbed into the constants; a full level structure does not give an arithmetic image literally equal to the unrestricted adelic group.\par
\section*{Checks of the principal mechanism}
\subsection*{Finite jets and monodromy}
The tame-ramification argument in Theorem 2.1 uses semistable reduction and unipotence of prime-to-\(p\) torsion. Lemma 3.2 is a finite-group argument: a subgroup meeting every conjugacy class is the whole finite group, and a rational point with trivial Frobenius image removes the arithmetic quotient. Its stated conclusion is modulo \(p^2\), not an automatic integral lifting result. Problem 1 isolates that distinction.\par
The simultaneous finite quotient \(H\times H'\) is important. After surjectivity to the separate infinite factors, Goursat’s lemma reduces entanglement to a common quotient. Factoring through \(H\) and \(H'\) then makes surjectivity onto their product force that quotient to be trivial. This extra hypothesis also supplies the missing characteristic-two argument below. The statement about density one of smooth hypersurface sections on page 9 is stronger than needed: at a fixed rational point of a smooth \(d\)-fold, the asymptotic probability that a random section vanishes with all first derivatives is \(q^{-(d+1)}>0\). The required existence should instead use Bertini with the finite local conditions imposed and positive density of smooth transverse sections. No failure of the required curve is inferred from that quantitative wording.\par
\subsection*{Uniform preservation of the required formal jets}
Proposition 4.1 has the correct matrix orientation \[Q'=(X^{T})^{-1}QY.\] Both the actual PDF on page 11 and the TeX have this formula. An apparent inverse-placement discrepancy in extracted text was rejected after inspection of the page. The lifting criterion in Katz, Theorem 2.1(4), is the compatibility \[q_A(e,\phi^\vee m')=q_B(\phi e,m'),\] which gives precisely \(QY=X^TQ'\) [katz]. Problem 2 concerns the definition of fractional multiplication, not this matrix orientation.\par
Equation reference: jet. Here is a complete version of the order calculation needed in Theorem 5.1. Choose the source parameter so that \(q'_{11}=1+t\) exactly. For \((i,j)\ne(1,1)\), the prescribed branch satisfies \(q'_{ij}=1+O(t^{m+1})\). For a Hecke type \(n\geq n_1\geq\cdots\geq n_g\geq0\), put \[c_{ij}=n-n_i-n_j,\qquad a=n-2n_1,\qquad
 d_{ij}=c_{ij}-a=2n_1-n_i-n_j\geq0.\] Over the perfect residue field let \(f^{[p^a]}\) mean coefficientwise Frobenius, allowing its inverse when \(a<0\). After the finite purely inseparable change needed to define the isogeny, and removal of a common inseparable parametrization, the reduced image has parameter \(u=t^{p^a}\) and coordinates \[
 R_{ij}(u)=\bigl(q_{ij}^{\prime\,[p^a]}(u)\bigr)^{p^{d_{ij}}}.\] This is obtained by substituting into \((q'_{ij}(t))^{p^{c_{ij}}}\). In particular, \[R_{11}=1+u,\qquad
 R_{ij}=1+O(u^{(m+1)p^{d_{ij}}})=1+O(u^{m+1})
 \quad ((i,j)\ne(1,1)).\] The first equality makes the reduced image formally smooth. The inequalities are uniform in the Hecke degree and valid for \(p=2\). Thus the requested primitive \(m\)-jet survives. The printed parameter need not itself be a uniformizer of the reduced image; the distinction is essential to this verification.\par
The hypersymmetric argument also has the required scope. An isogeny over \(\mathbb F_q\) to \(E^g\) makes Frobenius scalar on the rational Tate module and hence on its stable integral lattice. It preserves every finite \(p\)-subgroup and every rank-one cocharacter submodule. The resulting formal subtorus may be nonsplit over \(\mathbb F_q\); scalar descent, not split coordinates over \(\mathbb F_q\), is what is used.\par
\subsection*{Density and the uniform jet bound}
The proof of Proposition 5.3 is read in centered coordinates \(z_{ij}=q_{ij}-1\), with \(i\leq j\). A vector with all coordinates divisible by \(p\) does not itself generate a saturated line, but its parametrized image lies in the formal torus associated with the saturation. This removes the apparent primitivity obstruction in the chosen test family.\par
For completeness, a nonzero formal series has finitely many componentwise minimal exponent vectors, by Dickson’s lemma. For each such vector the substitution in the proof has weight \[W_N(\alpha)=\sum_{i\leq j}\alpha_{ij}p^{N^i+N^j}.\] For two distinct minimal vectors, the largest pair in the order determined by \(N^i+N^j\) eventually dominates their weight difference. There are finitely many pairs of minimal vectors, so for all sufficiently large \(N\) exactly one of them has least weight. Every other exponent dominates a minimal exponent and has strictly larger weight unless it is that exponent itself. Thus the substituted series cannot vanish. This argument includes characteristic two and does not replace the formal power series by a polynomial.\par
The passage from pointwise density to Proposition 5.4 additionally uses algebraicity of the finite-jet containment locus. Conditional on the cited formal-torus structure of the diagonal completion, the finite Igusa-level calculation must be truncated to the stated order. The finitely many rank-one jet directions then give a closed containment locus. These loci descend in a Noetherian space; their empty intersection, supplied by Proposition 5.3, implies that one finite order suffices. The original Chai reference was inaccessible from the attempted official hosts. Its indexed transcription is compatible with the cited structure, but it is not recorded as an inspected primary PDF. This dependency is an explicit verification limit.\par
Together with geometric irreducibility of Hecke translates, these checks support the constant-divisor mechanism. For the moving-divisor step, Lemma 5.6 is a uniform finite intersection assertion. Its uncountable-field argument can be made precise by choosing successively compatible prefixes from nested constructible subsets of the curve; Noetherian stabilization then gives the contradiction. The degree of the alternating cover must also be chosen at least as large as the bound in that lemma. No unrestricted inference that every finite extension preserves monodromy is used; see Problem 3.\par
\section*{Numbered concerns}
Problem 1 (Reduction modulo four does not alone force full \(2\)-adic monodromy).  Assessment: confirmed Minor. The standalone implication in the proof of Theorem 3.1, PDF page 9, TeX line 468, is false at \(p=2\) for every \(g>1\).\par
Let \(G=\operatorname{GL}_g(\mathbb Z_2)\) and set \[J_0=\{M\in G:\det M\bmod8\in\{1,7\}\}.\] This is a proper closed subgroup: \(\operatorname{diag}(3,1,\ldots,1)\) is excluded. Nevertheless \(J_0\) surjects onto \(\operatorname{GL}_g(\mathbb Z/4)\). Given any lift of an invertible matrix modulo four, multiplication by \(\operatorname{diag}(5,1,\ldots,1)\) preserves its reduction modulo four and switches its determinant to the required residue when necessary. The independent exact enumeration for \(g=2\) gives \[|\operatorname{GL}_2(\mathbb Z/8)|=1536,\quad |J_0\bmod8|=768,\quad
 |\operatorname{im}(J_0\bmod4)|=|\operatorname{GL}_2(\mathbb Z/4)|=96.\] The calculation is saved in `root\_exact\_checks\_001.json`.\par
Counterevidence and complete repair from existing data. The constructed curve is also required to surject onto \(G/X\), where every element of \(X=X_{2,a}\) is a \(2^a\)-th power and \(a\geq1\). This is stronger than the hypothesis just refuted. Let \(J\) be the actual closed image, surjective modulo four. In \(G_8=\operatorname{GL}_g(\mathbb Z/8)\) the kernel of reduction to \(G_4\) is \(N=1+4M_g(\mathbb F_2)\). Squares of lifts of \(1+2E_{ij}\), for \(i\ne j\), give \(1+4E_{ij}\). Commutators of lifts of \(1+2E_{ij}\) and \(1+2E_{ji}\) give \(1+4(E_{ii}+E_{jj})\). Consequently \(J_8\cap N\) contains every trace-zero matrix. If \(J_8\) were proper, it would have index two in \(G_8\) and would contain all squares. It would therefore contain the image of \(X\). But \(J\to G/X\) is surjective, so \(JX=G\), contradicting properness modulo eight. Thus \(J\) surjects modulo eight. Repeated squaring of congruence lifts then gives surjectivity at every higher power of two, hence \(J=G\). For odd \(p\), the usual congruence-power argument starts from modulo \(p^2\).\par
The defect is local because the hypotheses already imposed in the paper prove the needed conclusion. The subgroup \(J_0\) is not presented as a geometric counterexample satisfying all those hypotheses, and Theorem 3.1 is not refuted.\par
Problem 2 (Inverse Frobenius is missing from fractional multiplication).  Assessment: confirmed Minor. Section 4.1, PDF page 10, TeX line 496, defines \(p^{-n}Q\) by replacing \(t\) with \(t^{1/p^n}\) while leaving coefficients unchanged. This does not extend the natural \(\mathbb Z_p\)-module structure on \(1+t k[[t]]\) to inverse \(p\)-multiplication.\par
Take \(p=2\), \(k=\overline{\mathbb F}_2\), and \(\alpha\in\mathbb F_4\) with \(\alpha^2+\alpha+1=0\). For the admissible multiplicative coordinate \(q=1+\alpha t\), the printed prescription gives \[q_{\rm printed}^{1/2}=1+\alpha t^{1/2},\qquad
  (q_{\rm printed}^{1/2})^2=1+(\alpha+1)t\ne q.\] The correct square root is \(1+(\alpha+1)t^{1/2}\). For \(g=2\), put this coordinate in \(q_{11}\) and set all other symmetric coordinates equal to one. This is a polarized ordinary deformation. With isogeny matrices \(X=2I_2\), \(Y=I_2\), the lifting criterion requires \((q'_{11})^2=q_{11}\). The printed scalar operation fails precisely that requirement; it is not a complaint about optional notation.\par
For general \(p\) the corrected operation is \[\left(\sum_{\nu\geq0}a_\nu t^\nu\right)^{1/p^n}
       =\sum_{\nu\geq0}a_\nu^{1/p^n}t^{\nu/p^n}.\] Parameter substitution agrees with it when all coefficients are fixed by the relevant Frobenius. That valid special case explains the examples with coefficient one, but not the stated arbitrary algebraically closed residue field.\par
The explicit finite-field calculation and the Katz lifting criterion establish the discrepancy. Its demonstrated impact is a local coordinate operation. Formula [jet] (\#jet) checks the entire order argument actually needed for Theorem 5.1 and shows that coefficient Frobenius preserves all its vanishing-order conditions. The core jet-preservation conclusion therefore survives for every type; the Minor assessment rests on that bounded impact, not merely on the shortness of the correction.\par
Problem 3 (The constituent condition in Lemma 5.5 needs a precise meaning).  Assessment: Needs Evidence / scope clarification. Lemma 5.5, PDF page 17, TeX lines 702–710, does not require \(\phi:\operatorname{Gal}_K\to H\) to be surjective and uses an undefined group-theoretic constituent condition.\par
Under the composition-factor reading, the assertion is false. Take \(K=\mathbb F_q\), \(L=\mathbb F_{q^2}\), \(G=C_2\), and \(H=A_5\). Map a Frobenius generator of \(\operatorname{Gal}_K=\widehat{\mathbb Z}\) to \((12)(34)\in A_5\). The groups \(C_2\) and \(A_5\) have no common simple composition factor, but \[\phi(\operatorname{Gal}_K)=C_2,\qquad \phi(\operatorname{Gal}_L)=1.\] Every field and continuity hypothesis holds. The proof incorrectly treats a quotient of a subgroup of \(H\) as if it had to be a quotient of \(H\).\par
There are two substantive counterreadings. If the constituent condition excludes the relevant simple sections of all subgroups of \(H\), the example is inadmissible. Alternatively, requiring \(\phi\) to be surjective makes the proof valid under the composition-factor condition: after passing to the normal closure, the image quotient is a common quotient of \(H\) and \(G\). The actual application starts with the surjective monodromy constructed for \(C_0\), so that repair supplies the necessary application. Because the terminology admits a stronger reading, this is not counted as an unqualified confirmed false lemma. The report requests an explicit definition or the surjectivity hypothesis, without claiming failure of Theorem 1.1.\par
Problem 4 (The degree argument must specify correspondence multiplicities).  Assessment: Needs Evidence. Lemma 6.1 and equations (6.1.1)–(6.1.2), PDF pages 19–21, use \(\deg\tau(C)=\deg(\tau)\deg(C)\) and count coincident Hecke targets by asserting that their image has distinct branches. These statements require a precise convention for \(\tau(C)\).\par
Here is an admissible counterexample to a universal reduced-image degree formula. Work over \(\overline{\mathbb Q}\) with \(g=2\), \(r=3\), and \(\ell=7\). Choose a smooth proper curve \(B\subset\mathcal A_2[3]\) with maximal geometric monodromy, as provided by the complete intersection and Lefschetz construction used in Section 6.1. The Satake boundary has codimension two, so the complete intersection can avoid it. The generic abelian surface has endomorphism ring \(\mathbb Z\).\par
Let \(T\) be the Lagrangian \(7\)-isogeny correspondence. Its degree is \[d=(7+1)(7^2+1)=400.\] Since \(7\equiv1\pmod3\), the dual correspondence has compatible level structure and \(T\) is self-dual. The \(400\) targets of a generic point of \(B\) are distinct: two equivalent polarized quotient isogenies would, after composition with a dual, produce an endomorphism of a variety with endomorphism ring \(\mathbb Z\), forcing their kernels to agree. Monodromy acts transitively on the Lagrangians. Thus \(C=T(B)_{\rm red}\) is an irreducible proper curve.\par
For a generic \(b\in B\), all its \(400\) neighbours lie on \(C\). Self-duality gives at least \(400\) incidences above \(b\) in the correspondence applied to \(C\). Consequently the cycle \[Z_T(C)=p_{2*}p_1^*[C]\] contains \(B\) with coefficient at least \(400\), whereas \(T(C)_{\rm red}\) contains \(B\) only once. The Hodge bundle \(H\) has positive degree on \(B\) and on every proper image component. The correct projection formula gives \[\deg_H Z_T(C)=400\deg_H C,\qquad
 \deg_H T(C)_{\rm red}
       \leq400\deg_H C-399\deg_H B<400\deg_H C.\] At a generic point of the component \(B\), the reduced image is smooth, despite the multiple incidences.\par
Why this does not confirm a Major defect. Section 2.2 also uses labelled quotient points, and the standard action of a correspondence on a cycle retains precisely these multiplicities. Under that natural reading, Lemma 6.1’s projection formula is correct. The example refutes the reduced reading; it does not refute the cycle reading or establish that the special curve used later has the offending generic degree. The earlier explicit reduced-scheme convention occurs in the characteristic-\(p\) discussion, so its unqualified transfer to Section 6.1 cannot simply be assumed. This admissible alternative interpretation prevents confirmation of a false core equality.\par
A complete counting formulation. Let \(Y_\tau\) be the normalized Hecke cover of the source curve, with maps \(f_1:Y_\tau\to C\) and \(f_2:Y_\tau\to\mathcal A_g[r]\). Count the divisor \(f_2^*V\) on this cover. The Hodge-bundle identity gives \[\deg f_2^*V=d_\tau(C\cdot V).\] For a primitive source point of degree \(n\) with monodromy index at most \(c\), one target in \(V\) forces at least \(d_\tau/c\) labelled incidences over each source conjugate. These are different points of the cover, even if their images coincide. Hence \[\deg f_2^*V\geq n\,d_\tau/c,\] contradicting \(n>c(C\cdot V)\). This formulation does not require the reduced image to be singular at a repeated target. It also corrects the separate claim on page 21 that \(\tau(x)\cap V\) itself is stable under \(\operatorname{Gal}_{\mathbb Q}\): it is the union over conjugates of \(x\) that is stable.\par
For Theorem 1.4, assuming \(\tau(x)\subset D\) counts the whole Hecke fibre over each Frobenius conjugate. Thus no large Galois orbit of an individual finite-field target is required. A proper compactification and the compatible Hecke/Hodge intersection class supply the same upper bound. These compactification extension properties were not independently re-established in this review. The report does not assert failure of Theorems 6.3 or 1.4; it asks that the cycle convention and compactification details be made explicit.\par
Problem 5 (Hilbert specialization does not supply the stated Galois field).  Assessment: Needs Evidence. Lemma 6.2, PDF pages 19–20, requires \(K/\mathbb Q\) to be Galois and \(x\) primitive over \(K\). Its proof produces conjugate degree-\(d\) fields \(K_i/L\) and then sets \(K=K_1\).\par
Irreducibility of a fibre does not imply normality of its residue field. For example \(X^3-2\) is irreducible over the Galois base \(L=\mathbb Q\), has discriminant \(-108\), and defines a non-Galois cubic field whose normal closure has group \(S_3\). This is a counterexample to the field-theoretic inference, not an asserted counterexample to the lemma with its moduli and monodromy hypotheses. Passing to a normal closure makes the extension Galois but does not make the original point primitive over the larger field. Neither condition can be silently exchanged for the other.\par
The cited Zywina Theorem 1.1 and Section 1.2 were checked in the acquired v1 PDF, pages 1–3 [zywina]. They support bounded specialization index and the dependence of the constant on the arithmetic and geometric monodromy groups. They do not add a normality conclusion for \(K_1\).\par
A weaker statement is sufficient for the intersection argument: arbitrarily large primitive degree with a uniform index bound, without normality. If \(y\) is chosen to generate the auxiliary field \(L\) over the fixed field of definition, then \(y=\pi(x_1)\) implies \(L\subseteq\mathbb Q(x_1)\) (with the fixed base inserted when needed). Thus \(K_1=L(x_1)\) is the minimal field of the point and has degree at least \([L:\mathbb Q]\). Counting all embeddings of \(K_1\) into \(\overline{\mathbb Q}\) replaces the use of \(\operatorname{Gal}(K_1/\mathbb Q)\) in Problem 4. This describes the necessary repair; the stronger Galois assertion itself remains unproved in the inspected argument. The construction of arbitrarily large auxiliary fields disjoint from the full monodromy field is a further structural input, not certified merely by the finite cubic example.\par
\section*{Overall impact and state}
Problems 1 and 2 establish local failures of stated algebraic operations or implications, with explicit counterexamples and full checks of their relevance. Existing finite-quotient data repairs the first; the uniform reduced-branch calculation establishes the jet consequence needed despite the second. The principal function-field existence theorem has not been refuted.\par
Problems 3–5 retain their uncertainty. In particular, the reduced-image example must not be relabelled as a counterexample to a correspondence-cycle identity. The number-field proof can be organized around a weaker specialization lemma and counting on the Hecke cover, but the stated stronger lemma requires its own justification. No separate correction or associated executable code resolving these items was identified in the bounded search. Absence of code is not a mathematical defect.\par
9 A. N. Shankar and J. Tsimerman, Abelian varieties not isogenous to Jacobians over global fields, arXiv:2105.02998v3 (5 March 2025). https://arxiv.org/abs/2105.02998v3. N. M. Katz, Serre–Tate local moduli, in Algebraic Surfaces (Orsay, 1976–78), Lecture Notes in Mathematics 868, Springer, 1981, 138–202. Consulted author-hosted scan, especially Theorem 2.1(4), printed pages 148–149. D. Zywina, Families of abelian varieties and large Galois images, arXiv:1910.14174v1 (30 October 2019), especially Theorem 1.1 and Section 1.2. https://arxiv.org/abs/1910.14174v1.\par
\section*{Supporting checks}
The thirteen current manifest inputs total 1,458,354 bytes. Their current sizes and SHA-256 values match the frozen manifest. The reviewed PDF is 322,644 bytes, SHA-256\par
`36f1e4ad23a65a2e0d2031997433bcbfbb77ab3edfe26c9cf5bd5d38ac294a94`.\par
The manifest SHA-256 is\par
`491ed6ba588cec0525105a1ef9b4417cda49b4c873f58b1b581cdac4486c5afc`.\par
The complete twenty-two source-page images were opened and inspected by the primary executor. The source archive is a gzip-compressed single TeX payload despite its stored suffix; it was decoded for reading, not compiled or executed.\par
Exact reviewer-authored calculations record the \(1536/768/96\) matrix counts and the \(\mathbb F_4\) coefficient discrepancy. The native monodromy notes preserve the rejected column-vector coset objection; a consistent row-vector convention removes it. The source/PDF inverse-placement suspicion was also rejected. The degree/cycle analysis preserves the bounded helper’s counterinterpretation and the resulting decision not to confirm a Major failure. Each problem has a separate evidence record.\par
The primary executor performed the scientific synthesis, report writing, compilation and actual report-page inspection. Three native bounded helpers, with the same host self-description, checked monodromy, Serre–Tate calculations, and field-degree/intersection arguments respectively. They are AI cross-checks, not human reviews or independent acceptance.\par
The current TeX, compiler, dependencies, cache changes, logs, PDF and every inspected image are bound by hashes in the current build and provenance records. Render receipts retain their original false inspection flags; actual page inspections are separately attested after viewing. Failed attempts and prior versions are retained. The original Chai PDF access failures are recorded, with no primary-page inspection claim for that dependency. The initial-review decision is pending; human reviews are empty.\par
\end{document}
